\documentclass[../thesis.tex]{subfiles}

\begin{document}
	\section{Обзор программной инфраструктуры}

	\subsection{Роль программных сервисов}

	Программные сервисы выполняют задачи компьютерного моделирования исследуемого физического комплекса.
	К основным задачам можно отнести:

	\begin{itemize}
		\item Автоматизированное управление режимами работы комплекса;

		\item Аналитическая обработка данных;

		\item Визуальная интерпретация объекта исследования.
	\end{itemize}

	В идеальных условиях, совокупность программных сервисов образует цифровой двойник исследуемого
	комплекса, позволяя проводить ряд экспериментов без непосредственного доступа к физическому оборудованию.

	Стоит обратить внимание, что изложенные выше задачи -- основные, но не исчерпывающие. Кроме них программными
	сервисами решаются и другие задачи. Например:

	\begin{itemize}
		\item Формирование базы знаний;

		\item Предоставление индивидуального пространства разработки.
	\end{itemize}

	\subsection{Категории пользователей и их функции}

	При взаимодействии с программным обеспечением лабораторного комплекса, пользователи, как правило,
	относятся к одной или нескольким ролям:

	\begin{itemize}
		\item Разработчик;

		\item Администратор;

		\item Оператор.
	\end{itemize}

	Каждая роль определяет набор сценариев взаимодействия с программным обеспечением.

	\subsection{Жизненный цикл программных сервисов}

	На протяжении всего периоды эксплуатации, программные сервисы проходят несколько этапов. Типовой жизненный
	цикл программного сервиса приведен на рисунке \ref{fig:1}.

	\begin{figure}[H]
		\centering
		\includegraphics[width=1\textwidth]{assets/img/p1/lifecycle.png}
		\caption{Жизненный цикл программных сервисов}
		\label{fig:1}
	\end{figure}

	\subsection{Особенности}

	Основная особенность заключается в вычислительных и административных ресурсах лаборатории:

	\begin{itemize}
		\item Небольшая команда и ограниченные административные ресурсы;

		\item Локальная инфраструктура. Ограниченное количество вычислительных узлов расположенных в
			непосредственной близости к объекту исследования;

		\item Гетерогенные сервисы;

		\item Длительный срок эксплуатации сервисов. Некоторые сервисы предполагают непрерывный
			режим работы в течение длительного периода времени (год и более);

		\item Большой объем существующих разработок, требующих обеспечения обратной совместимости.
	\end{itemize}

	\section{Анализ существующей модели эксплуатации}

	Анализ проблем целесообразно проводить по категориям, рассматривая и решая задачи по порядку:

	\begin{itemize}
		\item Управление исходным кодом и среда разработки;

		\item Автоматизация сборки, тестирования и развертывания;

		\item Конфигурация;

		\item Координация времени выполнения;

		\item Межсерверное взаимодействие;

		\item Графический интерфейс;

		\item Документация и поиск информации.
	\end{itemize}

	\noindent
	Рассмотрим каждый класс задач отдельно, отвечая на вопросы:

	\begin{itemize}
		\item Зачем он нужен?

		\item Как сделано сейчас?

		\item В чем проблема?

		\item Какие существуют подходы?
	\end{itemize}

	\subsection{Управление исходным кодом и среда разработки}

	В условиях небольшой команды и ограниченных административных ресурсов грамотное управление
	исходным кодом и обеспечение возможности совместной разработки необходимы для эффективного
	распределения задач, контроля версий, предотвращения конфликтов при слиянии кода и появления
	дублирующегося функционала, а также для сосуществования гетерогенных сервисов на разных языках и платформах.
	Грамотный подход позволяет минимизировать риски ошибок, ускорить разработку, повысить качество
	разрабатываемого кода и обеспечить долгосрочную поддерживаемость сервисов.

	В существующей модели каждый сервис сопровождается независимо от остальных. Отсутствует стандартный
	подход к управлению внутренними зависимостями. Среда разработки не зафиксирована: разработчику необходимо
	найти и установить нужные зависимости и инструменты самостоятельно, попутно решая возникающие
	конфликты.

	С таким подходом сопряжено несколько проблем:

	\begin{itemize}
		\item Большой объем дублирующегося функционала;

		\item Высокий порог входа для разработчиков;

		\item Отсутствие воспроизводимости среды разработки.
	\end{itemize}

	Современные практики предполагают стандартизацию среды разработки посредством ее контейнеризации (см.
	\cite{devcontainers}). Помимо решения проблемы воспроизводимости и повышению удобства входа,
	такой подход открывает новые возможности для создания удаленных окружений разработки (RDE, см. \cite{rde},
	\cite{rde-uber}, \cite{tailscale}).

	Для управления зависимостями в гетерогенной среде существует два подхода:

	\begin{enumerate}
		\item Поддержка инструментов управления артефактами;

		\item Коллокация исходного кода (см. \cite{monorepo}, \cite{monorepo-google}).
	\end{enumerate}

	\subsection{Автоматизация сборки, тестирования и развертывания}

	Автоматизация сборки, тестирования и развертывания критична для обеспечения надежности и скорости
	доставки сервисов. В условиях гетерогенных сервисов на разных языках и платформах, автоматизация снижает
	вероятность ошибок, обеспечивает воспроизводимость процессов и позволяет минимизировать ручной
	труд небольшой команды. Грамотно выстроенный процесс интеграции значительно ускоряет процесс проверки
	новых концепций, что критично в экспериментальных лабораториях.

	В рассматриваемой модели все инструменты, включая систему сборки, изолируются одним проектом. Стандартной
	системы управления внешними зависимостями нет. Большинство проектов полагаются на наличие
	системных пакетов. Процесс тестирования осуществляется при непосредственном взаимодействии с рабочей
	средой. Развертывание сервисов осуществляется администратором в ручном режиме.

	Сопутствующие такому подходу проблемы включают:

	\begin{itemize}
		\item Невозможность разработки и использования общих внутренних инструментов;

		\item Высокий порог входа для администратора;

		\item Невоспроизводимость и риск конфликтов с рабочей средой;

		\item Длительный процесс интеграционного тестирования;

		\item Высокая нагрузка на администратора и длительный процесс ввода в эксплуатацию.
	\end{itemize}

	На практике подобные проблемы решаются следующими подходами:

	\begin{itemize}
		\item Коллокации исходного кода и внедрения общих инструментов;

		\item Использованию инструментов управления внешними зависимостями;

		\item Созданием изолированной тестовой среды;

		\item Использованию оркестраторов и систем управления инфраструктурой.
	\end{itemize}

	\subsection{Координация времени выполнения}

	Подумать надо.

	\subsection{Конфигурация}

	Конфигурации сервисов в лабораторной среде являются средством и результатом любой
	экспериментальной деятельности. Состояние всего комплекса описывается набором конфигураций, а
	значит грамотное управление последними является первостепенной задачей.

	В существующей модели для конфигурации сервисов применяются различные форматы, определяемые разработчиком.
	Располагается конфигурация в файловой системе ВМ, на которой запускается сам сервис. Некоторые
	конфигурации располагаются на сетевом диске. Оператор, как правило, взаимодействует с конфигурацией
	посредством графического интерфейса, поставляемого вместе с сервисом. Некоторые сервисы не предусматривают
	наличие внешней конфигурации и хранят состояние только в памяти.

	Подобный подход включает следующие очевидные недостатки:

	\begin{itemize}
		\item Высокая нагрузка на администратора по управлению конфигурациями;

		\item Отсутствие версионирования;

		\item Незащищенность от несанкционированного доступа;

		\item Отсутствие унифицированного механизма конфигурирования сервисов.
	\end{itemize}

	\subsection{Межсервисное взаимодействие}

	Система межсервисного взаимодействия определяет возможности обработки данных и применения различных
	вычислительных моделей, упрощающих разработку сложных алгоритмов (см. \cite{mapreduce}, \cite{dataflow}).
	Кроме этого, общая производительность системы не может превзойти уровень, допускаемый общей шиной
	данных.

	Преобладающей моделью взаимодействия в лабораторных комплексах является событийно-ориентированная
	архитектура. Основной информационный поток активно генерируется устройствами, непосредственно взаимодействующими
	с объектом исследования, собирая информацию о состоянии подсистем и записывая данные в один или
	несколько информационных каналов. Дополнительным источником информации служат события,
	поступающие со стороны оператора, например команда о переводе установки из одного режима в другой.

	В существующей модели, каналы данных обслуживаются собственным брокером сообщений с внутренним
	текстовым протоколом. Формального описания типа передаваемого сообщения нет.

	Возникающие проблемы включают:

	\begin{itemize}
		\item Снижение производительности за счет сериализации строк;

		\item Высокая связность сервисов и риск каскадных сбоев;

		\item Отсутствие возможности эволюции типа полезной нагрузки.
	\end{itemize}

	На практике взаимодействие между сервисами разделяется на два уровня. На первом уровне происходит
	обмен информацией с устройствами, которые взаимодействуют непосредственно изучаемым объектом. На втором
	уровне происходит обмен информацией между аналитическими сервисами, выполняющими основной объем
	вычислений.

	На первом уровне данные передаются в бинарном формате, обеспечивающем быструю сериализацию и строгую
	типизацию.

	На втором уровне данные сериализуются в эффективный по хранению
	\linebreak
	структурированный формат, сопровождаемый некоторой схемой.

	\subsection{Графический интерфейс}

	В лабораторной эксплуатации графический интерфейс это основной метод управления и наблюдения.
	Качество предоставляемых интерфейсов напрямую влияет на работу оператора и результаты
	экспериментальной деятельности.

	В существующей модели графический интерфейс интегрируется с сервис, который поставляется в форме стационарного
	приложения. Хотя унифицированного содержания интерфейса нет, большинство сервисов включают инструменты
	для конфигурирования, запуска, остановки, наблюдения состояния и визуализации результатов работы.

	Проблемы и недостатки, сопряженные с таким подходом:

	\begin{itemize}
		\item Зависимость от инструментов с возможностью создания графического интерфейса;

		\item Отсутствие многопользовательского и удаленного доступа к сервису;

		\item Ограничение среды выполнения сервиса ВМ оператора;

		\item Дополнительные временные затраты на реализацию типовых графических элементов.
	\end{itemize}

	На практике, с целью минимизации времени разработки, повышению производительности и обеспечению
	многопользовательского доступа стараются
	\linebreak
	разделять вычислительную часть сервиса и графическую.

	Для конфигурации и управления сервисами применяются инструменты построения внутренних приложений
	с Web интерфейсом и функционалом графического конструктора.

	Для визуализации используют инструменты с Web интерфейсом, которые предоставляют большой набор конфигурируемых
	визуализаций под различные задачи с возможность реализации собственных под специфические нужды.

	\subsection{Документация и поиск информации}

	Документация -- результат работы комплекса. По важности в эксплуатации, качественная документация
	сравнима с корректной конфигурацией.

	В существующей модели документация к сервисам либо отсутствует, либо физически распределена.

	Такой подход ведет к потере знаний, существенному увеличению входного порога и постепенной утрате
	согласованности различных аспектов работы лабораторного комплекса.

	\subsection{Вывод}

	Единой системы, которая бы решала все эти классы задач нет, но каждую задачу можно решить в
	отдельности существующими средствами. Некоторые открытые системы обладают излишним в лабораторной
	эксплуатации функционалом, а некоторые недостаточным.

	\section{Требования к решению}

	\begin{enumerate}
		\item Воспроизводимость среды разработки и развертывания;

		\item Унифицированные процессы разработки, регистрации, конфигурации, документации, тестирования
			и развертывания сервисов;

		\item Согласованное, прозрачное состояние системы с единственным источником истины;

		\item Совместимость с существующими сервисами и протоколами;

		\item Формальная модель данных для межсервисных коммуникаций;

		\item Централизованный графический интерфейс оператора с возможностью запускать и останавливать
			сервисы, наблюдать их состояние, осуществлять конфигурацию, а также визуализировать
			результаты их работы;

		\item Документация должны быть доступна централизованно, с возможностью поиска информации;

		\item Модульность и масштабируемость;

		\item Низкая эксплуатационная и административная нагрузка;

		\item Быстрая адаптация разработчиков, операторов и администраторов;

		\item Возможность поэтапного внедрения без остановки действующих процессов.
	\end{enumerate}

	Ключевым требованием к предлагаемому решению является минимизация инфраструктурной нагрузки. Вследствие
	этого, приоритет нужно отдавать распространенным открытым компонентам с активной поддержкой
	сообщества, в то время как собственная разработка ограничивается прозрачным интеграционным слоем,
	необходимым для адаптации этих компонентов под специфические требования лаборатории.

	\section{Постановка задачи}

	Основываясь на проведенном анализе контекста лаборатории и существующих открытых решений, в
	рамках настоящей работы ставятся следующие задачи:

	\begin{enumerate}
		\item Разработка новой архитектуры, решающей обозначенные проблемы. Вывод критериев сравнения
			существующих решений на соответствие требованиям лаборатории;

		\item Разработка унифицированной модели представления гетерогенных сервисов;

		\item Разработка интеграционных слоев, необходимых для реализации новой архитектуры;

		\item Стандартизация и документация процессов;

		\item Построение интеграционного плана, обеспечивающего поэтапную модернизацию рабочего
			процесса.
	\end{enumerate}
\end{document}